%% heterogeneity.tex — Section 6: Heterogeneity
%%
%% [INSERT] markers requiring filled values:
%%   Table 8 (or inline): gender × May2021 interaction estimates (6.1)
%%   Table 8 (or inline): rural × May2021 interaction estimates (6.2)
%%   Table 8 (or inline): SES tertile × May2021 interaction estimates (6.3)
%%   Table 9: remote learning duration estimates among 2021 students (6.4)
%%   [INSERT: interaction coefficient] throughout — replace with actual estimates
%%
%% Equation (2) is the augmented interaction specification.
%% Equation (3) is the within-2021 remote-duration specification.

The main estimates in Section~\ref{sec:results} identify the average
cross-cohort gap for the full sample. I now examine whether that gap
varies systematically across student subgroups. For binary characteristics,
I augment Equation~(1) with an interaction term:

\begin{equation}
Y_{irt} = \alpha_r + \beta_1 \cdot \mathit{May2021}_t
          + \beta_2 \cdot Z_i
          + \beta_3 \cdot (\mathit{May2021}_t \times Z_i)
          + X'_{irt} \beta_4 + \varepsilon_{irt},
\label{eq:interaction}
\end{equation}

\noindent where $Z_i$ is the characteristic of interest, $X_{irt}$
retains the full control set of Equation~(1), and $\beta_3$ identifies
whether the cross-cohort gap differs for students with $Z_i = 1$
relative to those with $Z_i = 0$. Region fixed effects $\alpha_r$ and
school-level clustered standard errors are maintained throughout.
The main effect $\beta_1$ now represents the gap for the reference group
($Z_i = 0$), and $\beta_1 + \beta_3$ is the gap for students with
$Z_i = 1$. All interactions are estimated separately for mathematics and
reading.

\subsection{Gender}
\label{subsec:gender}

Table~5 shows a large female advantage in reading of 0.288~SD, a pattern
well documented in early literacy research \citep{fryer2010empirical}.
In mathematics, the baseline gender gap is small and statistically
indistinguishable from zero ($-0.012$~SD for girls relative to boys).
I test whether the pandemic widened or narrowed these gender gaps by
estimating Equation~(\ref{eq:interaction}) with $Z_i =
\mathit{Female}_i$.

The interaction estimates appear in Table~8. For
mathematics, $\hat{\beta}_3 =$ XX.XX~SD (SE = XX.XX, XX). %% PLACEHOLDER
For reading, $\hat{\beta}_3 =$ XX.XX~SD (SE = XX.XX, XX). %% PLACEHOLDER
XX.XX. %% PLACEHOLDER: 1-2 sentence interpretation of gender result

The international literature on gender heterogeneity in pandemic learning
losses is mixed. \citet{betthaeuser2023systematic} report no consistent
pattern across high-income countries. Evidence from low- and
middle-income countries is thinner, but \citet{patrinos2022analysis}
document similarly modest gender differences in their multi-country
analysis. The Ukraine estimates are XX.XX %% PLACEHOLDER: consistent/inconsistent
with this pattern.

\subsection{Rural/Urban Location}
\label{subsec:rural}

Rural students in Ukraine perform substantially below their urban peers:
Table~5 reports a gap of $-0.379$~SD in mathematics and $-0.290$~SD in
reading after full controls. Remote instruction during the pandemic
plausibly widened this gap further. Rural schools in Ukraine face
persistent disadvantages in digital infrastructure and household internet
access, both of which are necessary conditions for effective distance
learning. Ukraine's national broadband statistics show substantially lower
coverage rates outside oblast centres, and the teacher questionnaire
responses confirm that remote learning was less uniformly implemented
in rural schools.

I estimate Equation~(\ref{eq:interaction}) with $Z_i = \mathit{Rural}_i$.
The interaction $\hat{\beta}_3$ measures whether rural students
experienced a larger cross-cohort decline than urban students. The
estimates are reported in Table~8.

For mathematics, $\hat{\beta}_3 =$ XX.XX~SD (SE = XX.XX, XX). %% PLACEHOLDER
For reading, $\hat{\beta}_3 =$ XX.XX~SD (SE = XX.XX, XX). %% PLACEHOLDER
XX.XX. %% PLACEHOLDER: 1-2 sentence interpretation of rural result

The rural--urban dimension deserves particular attention in the Ukrainian
context. Prior to the pandemic, rural schools already contended with
teacher shortages and larger class sizes in remote oblasts
\citep[see][]{world_bank2021ukraine}. A significant interaction would
imply that pandemic losses compounded pre-existing structural
inequalities, with consequences for remediation targeting.

\subsection{Socioeconomic Status}
\label{subsec:ses}

A central finding in the international literature is that pandemic
learning losses were disproportionately large for disadvantaged students.
\citet{betthaeuser2023systematic} estimate that losses in studies
stratified by SES were on average larger for low-SES students by
approximately 0.05~SD across 42 studies. \citet{engzell2021learning}
document a gradient in the Netherlands whereby students from lower-educated
households experienced steeper losses during the lockdown period.

To assess the SES gradient in Ukraine, I extend Equation~(\ref{eq:interaction})
to include separate interactions for the SES-middle and SES-high tertile
dummies interacted with $\mathit{May2021}_t$, with SES-low as the omitted
reference group. The estimated interactions thus represent whether the
cross-cohort gap was smaller or larger for students in higher SES tertiles
relative to the lowest tertile.

Table~8 reports the SES interaction estimates. For mathematics:
\begin{itemize}
  \item SES-middle $\times$ May2021: $\hat{\beta}_3 =$ XX.XX~SD
        (SE = XX.XX, XX) %% PLACEHOLDER
  \item SES-high $\times$ May2021: $\hat{\beta}_3 =$ XX.XX~SD
        (SE = XX.XX, XX) %% PLACEHOLDER
\end{itemize}
For reading:
\begin{itemize}
  \item SES-middle $\times$ May2021: $\hat{\beta}_3 =$ XX.XX~SD
        (SE = XX.XX, XX) %% PLACEHOLDER
  \item SES-high $\times$ May2021: $\hat{\beta}_3 =$ XX.XX~SD
        (SE = XX.XX, XX) %% PLACEHOLDER
\end{itemize}
XX.XX. %% PLACEHOLDER: 1-2 sentence interpretation of SES gradient result

One important caveat applies to this analysis. As documented in
Section~\ref{sec:data}, the SES composition of the 2021 sample shifted
substantially toward higher tertiles relative to 2018. If low-SES
students disproportionately exited the 2021 sample --- precisely because
they were most affected by the disruption --- then the observed interaction
coefficients underestimate the true SES gradient in learning losses.
The reweighting exercise in Section~\ref{subsec:ipw} partially addresses
this concern but cannot fully recover the experiences of students who
were not tested.

\subsection{Remote Learning Duration}
\label{subsec:remote}

A distinctive feature of the 2021 wave is that the teacher questionnaire
collected information on the duration of remote instruction experienced
by each class during the 2019--20 and 2020--21 academic years. This
allows a within-2021 analysis of whether longer exposure to distance
learning is associated with lower test scores, conditional on all other
observables.

The distribution of remote learning exposure is as follows: 9.9 percent
of students were in classes that reported zero months of remote
instruction; 29.2 percent experienced less than one month; 48.2 percent
experienced 1--2 months; 10.3 percent experienced 3--4 months; and
2.4 percent experienced five or more months. The modal category --- nearly
half of all students --- reflects the national pattern of reopening in
autumn 2020 after the spring 2020 closure, with adaptive quarantine
protocols in most regions.

Because remote learning duration is observed only in the 2021 wave, I
estimate its association with achievement among 2021 students only.
Let $R_{ic}$ denote the remote learning duration category for student
$i$ in class $c$, and define indicator variables for each category with
the ``less than one month'' group as the reference. The estimating
equation is:

\begin{equation}
\begin{split}
Y_{irt} &= \gamma_0 + \gamma_1 \cdot \mathbf{1}[R_{ic} = 0]
           + \gamma_2 \cdot \mathbf{1}[R_{ic} \in \{1{-}2\}] \\
        &\quad + \gamma_3 \cdot \mathbf{1}[R_{ic} \in \{3{-}4\}]
           + \gamma_4 \cdot \mathbf{1}[R_{ic} \geq 5]
           + X'_{irt} \delta + \alpha_r + \varepsilon_{irt},
\end{split}
\label{eq:remote}
\end{equation}

\noindent where $\mathbf{1}[\cdot]$ denotes indicator functions for
duration categories and $X_{irt}$, $\alpha_r$ retain their definitions
from Equation~(1). Standard errors are clustered at the school level.
Note that Equation~(\ref{eq:remote}) is estimated in the 2021 sample
only; it does not identify the effect of remote learning relative to
in-person instruction in the counterfactual sense, but rather the
conditional association between remote duration and achievement among
students who experienced the pandemic.

The estimates appear in Table~9. Students in classes reporting 3--4
months of remote instruction scored XX.XX~SD lower in mathematics than
those in classes with less than one month of exposure (SE = XX.XX,
$p < $ XX.XX). %% PLACEHOLDER
The association in reading is XX.XX~SD (SE = XX.XX). %% PLACEHOLDER
Students in classes with five or more months of remote instruction ---
a small share of the sample --- show the largest deficit: XX.XX~SD in
mathematics and XX.XX~SD in reading. %% PLACEHOLDER
Classes reporting zero months of remote learning show XX.XX %% PLACEHOLDER
differences relative to the reference group, a pattern that could
reflect measurement error in the teacher reports or genuine
heterogeneity in how schools defined remote instruction.

This within-2021 association should be interpreted with caution.
Remote learning duration is endogenous: schools in regions with worse
infrastructure, higher local COVID-19 incidence, or weaker
administrative capacity were more likely to keep students at home
for longer, and these factors are also predictive of lower baseline
achievement. Region fixed effects absorb between-region variation in
infrastructure and COVID severity, but within-region sorting of students
across schools with different remote durations remains a concern.
The estimates in Table~9 are therefore best understood as descriptive
associations that motivate future research with richer data, rather
than as causal estimates of a remote-learning treatment effect.
